\paragraph{Statement E.20 (a piecewise linear Higman exclusion).}
Let $\operatorname{BPL}_+(\mathbb R)$ be the group of
orientation-preserving piecewise linear homeomorphisms of the line
whose breakpoints are locally finite and which are bi-Lipschitz.
The bi-Lipschitz constant may depend on the element.
Every homomorphism from the Higman group $H$ to this group is
trivial. Consequently, every group normally generated by images
of copies of $H$, including the group in A.28, has no nontrivial homomorphism
to any subgroup of $\operatorname{BPL}_+(\mathbb R)$.

\emph{Proof.}
First suppose that $uvu^{-1}=v^2$ in this group and that $v\ne1$
has a fixed point. A nonempty component of its support has a finite
endpoint $p$. On the side pointing into that component, local
finiteness of the breakpoints makes $v$ linear near $p$, with
one-sided slope $\lambda\ne1$. Indeed, $v(p)=p$, so slope one
would make that interval pointwise fixed. The map $u$ preserves
$\operatorname{Fix}(v)$ and its complementary components, because
$\operatorname{Fix}(v^2)=\operatorname{Fix}(v)$.
At every fixed point, the one-sided derivatives of $u$ and $u^{-1}$
cancel in differentiating $uvu^{-1}$. Orientation preserves the
chosen side. The relation therefore gives
\[
 v'_{\pm}(up)^2=v'_{\pm}(p),\qquad
 v'_{\pm}(u^{-k}p)=\lambda^{2^k}\quad(k\ge0).
\]
For $\lambda>1$ these slopes are unbounded, and for $0<\lambda<1$
they tend to zero. Both alternatives contradict the finite upper
and positive lower slope bounds supplied by the bi-Lipschitz
condition on the single element $v$. Thus a nontrivial $v$ in such
a relation has no fixed point.

Now suppose that
\[
 aba^{-1}=b^2,\qquad cac^{-1}=a^2.
\]
If $b\ne1$, it is fixed-point-free. An increasing topological change
of coordinates takes $b$ to $t\mapsto t+1$ or $t\mapsto t-1$:
choose a homeomorphism on a fundamental interval and extend using
the iterates of $b$. In either case the first relation gives
$a(t+1)=a(t)+2$ in these coordinates. Hence
\[
 a(t+n)-(t+n)=a(t)-t+n\qquad(n\in\mathbb Z).
\]
Continuity supplies a fixed point of $a$. Fixed points do not depend
on the chosen coordinates. Applying the preceding argument to the
second relation in the original coordinates gives $a=1$, and then
$b=b^2$ gives $b=1$, a contradiction. Every Higman generator is the
$b$-term of two consecutive such relations, after replacing all
named generators by their inverses if using the presentation in
E.4. Thus every generator image is trivial. Normal generation gives
the final assertion. \hfill$\square$

The local finiteness and global bi-Lipschitz hypotheses are part
of the assertion. In particular, this is not an exclusion of
arbitrary groups of orientation-preserving homeomorphisms of the
line. A finite global set of positive slopes is sufficient for the
bi-Lipschitz condition. The proof uses two consecutive Higman
relations; a single power-conjugacy relation can hold nontrivially
for affine maps of the line.

\paragraph{Statement E.21 (a countable-quotient block with one homology cost).}
Let $P$ be a perfect group with an actual two-dimensional
classifying space $X$, and let $q:P\twoheadrightarrow E$ be a
surjection onto a countable group. Choose $h\in P$ with $q(h)\ne1$.
There is an injective homomorphism $P\hookrightarrow G$, realized
by a subcomplex inclusion between actual two-dimensional
classifying spaces, such that
\[
 H_1(G;\mathbb Z)=0,\qquad
 H_2(G;\mathbb Z)\cong H_2(P;\mathbb Z)\oplus\mathbb Z,
 \tag{E.21.1}
\]
and no homomorphism $G\to E$ extends $q$. The construction adds
only countably many cells to $X$, and $G$ is normally generated
by the image of $P$.

\emph{Proof.}
Choose $c_0,c_1,\ldots\in P$ whose $q$-images enumerate $E$;
repetitions are allowed. Use the convention $[x,y]=xyx^{-1}y^{-1}$.
Start from $P*\langle a\rangle$ and set
\[
 L=\left\langle P,a,y_0,y_1,\ldots\ \middle|\
 h=[a,y_0],\quad
 y_n=[a c_n^{-1},y_{n+1}]\ (n\ge0)\right\rangle.
 \tag{E.21.2}
\]
This presentation is an increasing union of injective cyclic HNN
extensions, as follows. The first new relation is
\[
 y_0 a^{-1}y_0^{-1}=a^{-1}h.
\]
Both associated cyclic subgroups are infinite: each displayed
endpoint has nonzero $a$-exponent in $P*\langle a\rangle$.
After $y_n$ has been added, write $k_n=a c_n^{-1}$. The next
relation is
\[
 y_{n+1}k_n^{-1}y_{n+1}^{-1}=k_n^{-1}y_n.
\]
The left endpoint is nontrivial of infinite order in the embedded
initial free product. The right endpoint has exponent one in the
most recently added stable letter $y_n$. The homomorphism from
that finite HNN stage to $\mathbb Z$ which sends $y_n$ to one and
its base group to zero proves that this endpoint also has infinite
order. The HNN normal form theorem therefore embeds each stage
in the next one.

Use $X\vee S^1$ for the initial vertex space and attach a cylinder
on a circle for each cyclic HNN extension. All vertex spaces are
aspherical and all edge maps induce injections, so the usual tree
of universal covers proves asphericity at each finite stage.
These are actual two-dimensional CW spaces, with the previous
stage as a subcomplex. A sphere in the union has image in a finite
subcomplex and hence in some finite stage; it is null-homotopic
there. The union is therefore an actual two-dimensional $K(L,1)$
containing $X$.

For ordinary relative cellular chains over $X$, there is one new
one-cell $a$ and the one-cells $y_n$. The root two-cell
$h=[a,y_0]$ has zero relative boundary. Each subsequent two-cell
$y_n=[a c_n^{-1},y_{n+1}]$ has relative boundary $y_n$, up to the
choice of orientation: all occurrences of $a$ and $y_{n+1}$ have
cancelling exponent sums. Thus
\[
 H_2(K(L,1),X;\mathbb Z)=\mathbb Z,\qquad
 H_1(K(L,1),X;\mathbb Z)=\mathbb Z\langle a\rangle.
\]
The connecting homomorphism from this relative $H_2$ sends its
generator to the class of $h$ in $H_1(P;\mathbb Z)$, which is
zero by perfection. Since there are no relative three-cells, the
long exact sequence gives
\[
 H_2(L;\mathbb Z)\cong H_2(P;\mathbb Z)\oplus\mathbb Z,
 \qquad H_1(L;\mathbb Z)=\mathbb Z\langle a\rangle.
\]
The splitting is not claimed to be canonical.

To kill the remaining first homology while retaining normal
generation, first adjoin a stable letter $t_0$ with
\[
 t_0 a t_0^{-1}=h.
\]
Both endpoints have infinite order: $a$ lies in the embedded
initial free product, and $h\ne1$ lies in the torsion-free group
$P$. Next adjoin $t_1,t_2,\ldots$ with
\[
 t_n=[h,t_{n+1}]\qquad(n\ge0).
\]
These again are injective cyclic HNN extensions, since the $n$th
relation is
\[
 t_{n+1}h^{-1}t_{n+1}^{-1}=h^{-1}t_n.
\]
The left endpoint has infinite order in the embedded $P$; the right
endpoint has exponent one in the latest stable letter $t_n$.
Let $G$ be the union. The same graph-of-spaces argument supplies
an actual two-dimensional classifying space including $X$.
The first new HNN edge kills the primitive first-homology class
$a$ and introduces $t_0$. Each later edge kills $t_n$ and introduces
$t_{n+1}$. The edge map on first homology is injective at each of
these steps, so Mayer--Vietoris leaves $H_2$ unchanged. Taking the
union kills every such first-homology class. This proves (E.21.1).
Equivalently, in the full ordinary relative cellular chain complex,
the added boundaries are $a$ and then $t_n$; together with the
previous $y_n$ boundaries, they make every new one-cell a boundary
and leave exactly the one root two-cycle.

If $f:G\to E$ extended $q$, choose $n$ with
$f(a)=q(c_n)$. Then $f(k_n)=1$, so the $n$th tail relation gives
$f(y_n)=1$. Successive earlier tail relations imply
$f(y_{n-1})=\cdots=f(y_0)=1$. The root relation now gives
$q(h)=[f(a),f(y_0)]=1$, a contradiction.

Finally, quotient by the normal closure of $P$. Then $h=1$,
so $t_0 a t_0^{-1}=h$ gives $a=1$. The relations for the $y_n$
now give $y_n=1$ for every $n$, and the relations
$t_n=[h,t_{n+1}]$ give $t_n=1$ for every $n$.
Thus the quotient is trivial, proving normal generation.
\hfill$\square$

Each use of this block adds an independent integral second-homology
class. In an increasing construction by these actual subcomplex
inclusions, subsequent two-dimensional attachments cannot remove
any existing cellular two-cycle. Thus the block proves a normally generating quotient
obstruction with an exact cost; it does not give a homology-preserving
replacement for A.27 or settle E2. If instead the root $h$ belongs
to a fresh Higman copy, the weaker conclusion is merely that this
fresh copy must map trivially whenever the image of $a$ lies in the
image of $P$. The old-coefficient HNN construction above is what
makes the prescribed extension $q$ itself impossible.


